\section{Pure ideals}
\label{section-pure-ideals}

\noindent
The material in this section is discussed in many papers, see for example
\cite{Lazard}, \cite{Bkouche}, and \cite{DeMarco}.

\begin{definition}
\label{definition-pure-ideal}
Let $R$ be a ring. We say that $I \subset R$ is {\it pure}
if the quotient ring $R/I$ is flat over $R$.
\end{definition}

\begin{lemma}
\label{lemma-pure}
Let $R$ be a ring.
Let $I \subset R$ be an ideal.
The following are equivalent:
\begin{enumerate}
\item $I$ is pure,
\item for every ideal $J \subset R$ we have $J \cap I = IJ$,
\item for every finitely generated ideal $J \subset R$ we have
$J \cap I = JI$,
\item for every $x \in R$ we have $(x) \cap I = xI$,
\item for every $x \in I$ we have $x = yx$ for some $y \in I$,
\item for every $x_1, \ldots, x_n \in I$ there exists a
$y \in I$ such that $x_i = yx_i$ for all $i = 1, \ldots, n$,
\item for every prime $\mathfrak p$ of $R$ we have
$IR_{\mathfrak p} = 0$ or $IR_{\mathfrak p} = R_{\mathfrak p}$,
\item $\text{Supp}(I) = \Spec(R) \setminus V(I)$,
\item $I$ is the kernel of the map $R \to (1 + I)^{-1}R$,
\item $R/I \cong S^{-1}R$ as $R$-algebras for some multiplicative
subset $S$ of $R$, and
\item $R/I \cong (1 + I)^{-1}R$ as $R$-algebras.
\end{enumerate}
\end{lemma}

\begin{proof}
For any ideal $J$ of $R$ we have the short exact sequence
$0 \to J \to R \to R/J \to 0$. Tensoring with $R/I$ we get
an exact sequence $J \otimes_R R/I \to R/I \to R/I + J \to 0$
and $J \otimes_R R/I = J/JI$. Thus the
equivalence of (1), (2), and (3) follows from
Lemma \ref{lemma-flat}. Moreover, these imply (4).

\medskip\noindent
The implication (4) $\Rightarrow$ (5) is trivial.
Assume (5) and let $x_1, \ldots, x_n \in I$.
Choose $y_i \in I$ such that $x_i = y_ix_i$.
Let $y \in I$ be the element such that
$1 - y = \prod_{i = 1, \ldots, n} (1 - y_i)$.
Then $x_i = yx_i$ for all $i = 1, \ldots, n$.
Hence (6) holds, and it follows that (5) $\Leftrightarrow$ (6).

\medskip\noindent
Assume (5). Let $x \in I$. Then $x = yx$ for some $y \in I$.
Hence $x(1 - y) = 0$, which shows that $x$ maps to zero in
$(1 + I)^{-1}R$. Of course the kernel of the map
$R \to (1 + I)^{-1}R$ is always contained in $I$. Hence we
see that (5) implies (9). Assume (9). Then for any $x \in I$
we see that $x(1 - y) = 0$ for some $y \in I$.
In other words, $x = yx$. We conclude that (5) is equivalent to (9).

\medskip\noindent
Assume (5). Let $\mathfrak p$ be a prime of $R$.
If $\mathfrak p \not \in V(I)$, then $IR_{\mathfrak p} = R_{\mathfrak p}$.
If $\mathfrak p \in V(I)$, in other words, if $I \subset \mathfrak p$,
then $x \in I$ implies $x(1 - y) = 0$ for some $y \in I$, implies
$x$ maps to zero in $R_{\mathfrak p}$, i.e., $IR_{\mathfrak p} = 0$.
Thus we see that (7) holds.

\medskip\noindent
Assume (7). Then $(R/I)_{\mathfrak p}$ is either $0$ or $R_{\mathfrak p}$
for any prime $\mathfrak p$ of $R$. Hence by
Lemma \ref{lemma-flat-localization}
we see that (1) holds. At this point we see that all of
(1) -- (7) and (9) are equivalent.

\medskip\noindent
As $IR_{\mathfrak p} = I_{\mathfrak p}$ we see that (7) implies (8).
Finally, if (8) holds, then this means exactly that $I_{\mathfrak p}$
is the zero module if and only if $\mathfrak p \in V(I)$, which
is clearly saying that (7) holds. Now (1) -- (9) are equivalent.

\medskip\noindent
Assume (1) -- (9) hold. Then $R/I \subset (1 + I)^{-1}R$ by (9) and
the map $R/I \to (1 + I)^{-1}R$ is also surjective by the description
of localizations at primes afforded by (7). Hence (11) holds.

\medskip\noindent
The implication (11) $\Rightarrow$ (10) is trivial.
And (10) implies that (1) holds because a localization of
$R$ is flat over $R$, see
Lemma \ref{lemma-flat-localization}.
\end{proof}

\begin{lemma}
\label{lemma-pure-ideal-determined-by-zero-set}
\begin{slogan}
Pure ideals are determined by their vanishing locus.
\end{slogan}
Let $R$ be a ring.
If $I, J \subset R$ are pure ideals, then $V(I) = V(J)$
implies $I = J$.
\end{lemma}

\begin{proof}
For example, by property (7) of
Lemma \ref{lemma-pure}
we see that
$I = \Ker(R \to \prod_{\mathfrak p \in V(I)} R_{\mathfrak p})$
can be recovered from the closed subset associated to it.
\end{proof}

\begin{lemma}
\label{lemma-pure-open-closed-specializations}
Let $R$ be a ring. The rule
$I \mapsto V(I)$
determines a bijection
$$
\{I \subset R \text{ pure}\}
\leftrightarrow
\{Z \subset \Spec(R)\text{ closed and closed under generalizations}\}
$$
\end{lemma}

\begin{proof}
Let $I$ be a pure ideal. Then since $R \to R/I$ is flat, by going down
generalizations lift along the map $\Spec(R/I) \to \Spec(R)$.
Hence $V(I)$ is closed under generalizations. This shows that the map
is well defined. By
Lemma \ref{lemma-pure-ideal-determined-by-zero-set}
the map is injective. Suppose that
$Z \subset \Spec(R)$ is closed and closed under generalizations.
Let $J \subset R$ be the radical ideal such that $Z = V(J)$.
Let $I = \{x \in R : x \in xJ\}$. Note that $I$ is an ideal:
if $x, y \in I$ then there exist $f, g \in J$ such that
$x = xf$ and $y = yg$. Then
$$
x + y = (x + y)(f + g - fg)
$$
Verification left to the reader.
We claim that $I$ is pure and that $V(I) = V(J)$.
If the claim is true then the map of the lemma is surjective and
the lemma holds.

\medskip\noindent
Note that $I \subset J$, so that $V(J) \subset V(I)$.
Let $I \subset \mathfrak p$ be a prime. Consider the multiplicative
subset $S = (R \setminus \mathfrak p)(1 + J)$. By definition of
$I$ and $I \subset \mathfrak p$ we see that $0 \not \in S$.
Hence we can find a prime $\mathfrak q$ of $R$ which is disjoint
from $S$, see
Lemmas \ref{lemma-localization-zero} and
\ref{lemma-spec-localization}.
Hence $\mathfrak q \subset \mathfrak p$ and
$\mathfrak q \cap (1 + J) = \emptyset$.
This implies that $\mathfrak q + J$ is a proper ideal of $R$.
Let $\mathfrak m$ be a maximal ideal containing $\mathfrak q + J$.
Then we get
$\mathfrak m \in V(J)$ and hence $\mathfrak q \in V(J) = Z$
as $Z$ was assumed to be closed under generalization.
This in turn implies $\mathfrak p \in V(J)$ as
$\mathfrak q \subset \mathfrak p$. Thus we see that $V(I) = V(J)$.

\medskip\noindent
Finally, since $V(I) = V(J)$ (and $J$ radical) we see that $J = \sqrt{I}$.
Pick $x \in I$, so that $x = xy$ for some $y \in J$ by definition.
Then $x = xy = xy^2 = \ldots = xy^n$. Since $y^n \in I$ for some $n > 0$
we conclude that property (5) of
Lemma \ref{lemma-pure}
holds and we see that $I$ is indeed pure.
\end{proof}

\begin{lemma}
\label{lemma-finitely-generated-pure-ideal}
Let $R$ be a ring. Let $I \subset R$ be an ideal.
The following are equivalent
\begin{enumerate}
\item $I$ is pure and finitely generated,
\item $I$ is generated by an idempotent,
\item $I$ is pure and $V(I)$ is open, and
\item $R/I$ is a projective $R$-module.
\end{enumerate}
\end{lemma}

\begin{proof}
If (1) holds, then $I = I \cap I = I^2$ by
Lemma \ref{lemma-pure}.
Hence $I$ is generated by an idempotent by
Lemma \ref{lemma-ideal-is-squared-union-connected}.
Thus (1) $\Rightarrow$ (2).
If (2) holds, then $I = (e)$ and $R = (1 - e) \oplus (e)$ as
an $R$-module hence $R/I$ is flat and $I$ is pure and $V(I) = D(1 - e)$
is open. Thus (2) $\Rightarrow$ (1) $+$ (3).
Finally, assume (3). Then $V(I)$ is open and closed, hence
$V(I) = D(1 - e)$ for some idempotent $e$ of $R$, see
Lemma \ref{lemma-disjoint-decomposition}.
The ideal $J = (e)$ is a pure ideal such that $V(J) = V(I)$ hence
$I = J$ by
Lemma \ref{lemma-pure-ideal-determined-by-zero-set}.
In this way we see that (3) $\Rightarrow$ (2). By
Lemma \ref{lemma-finite-projective}
we see that (4) is equivalent to the assertion that $I$ is pure
and $R/I$ finitely presented. Moreover, $R/I$ is finitely presented
if and only if $I$ is finitely generated, see Lemma \ref{lemma-extension}.
Hence (4) is equivalent to (1).
\end{proof}

\noindent
We can use the above to characterize those rings for which every finite
flat module is finitely presented.

\begin{lemma}
\label{lemma-finite-flat-module-finitely-presented}
Let $R$ be a ring. The following are equivalent:
\begin{enumerate}
\item every $Z \subset \Spec(R)$ which is closed and closed under
generalizations is also open, and
\item any finite flat $R$-module is finite locally free.
\end{enumerate}
\end{lemma}

\begin{proof}
If any finite flat $R$-module is finite locally free then the support
of $R/I$ where $I$ is a pure ideal is open. Hence the implication
(2) $\Rightarrow$ (1) follows from
Lemma \ref{lemma-pure-ideal-determined-by-zero-set}.

\medskip\noindent
For the converse assume that $R$ satisfies (1).
Let $M$ be a finite flat $R$-module.
The support $Z = \text{Supp}(M)$ of $M$ is closed, see
Lemma \ref{lemma-support-closed}.
On the other hand, if $\mathfrak p \subset \mathfrak p'$, then by
Lemma \ref{lemma-finite-flat-local}
the module $M_{\mathfrak p'}$ is free, and
$M_{\mathfrak p} = M_{\mathfrak p'} \otimes_{R_{\mathfrak p'}} R_{\mathfrak p}$
Hence
$\mathfrak p' \in \text{Supp}(M) \Rightarrow \mathfrak p \in \text{Supp}(M)$,
in other words, the support is closed under generalization.
As $R$ satisfies (1) we see that the support of $M$ is open and closed.
Suppose that $M$ is generated by $r$ elements $m_1, \ldots, m_r$.
The modules $\wedge^i(M)$, $i = 1, \ldots, r$ are finite flat $R$-modules
also, because $\wedge^i(M)_{\mathfrak p} = \wedge^i(M_{\mathfrak p})$
is free over $R_{\mathfrak p}$. Note that
$\text{Supp}(\wedge^{i + 1}(M)) \subset \text{Supp}(\wedge^i(M))$.
Thus we see that there exists a decomposition
$$
\Spec(R) = U_0 \amalg U_1 \amalg \ldots \amalg U_r
$$
by open and closed subsets such that the support of
$\wedge^i(M)$ is $U_r \cup \ldots \cup U_i$ for all $i = 0, \ldots, r$.
Let $\mathfrak p$ be a prime of $R$, and say $\mathfrak p \in U_i$.
Note that
$\wedge^i(M) \otimes_R \kappa(\mathfrak p) =
\wedge^i(M \otimes_R \kappa(\mathfrak p))$.
Hence, after possibly renumbering $m_1, \ldots, m_r$ we may assume that
$m_1, \ldots, m_i$ generate $M \otimes_R \kappa(\mathfrak p)$. By
Nakayama's Lemma \ref{lemma-NAK}
we get a surjection
$$
R_f^{\oplus i} \longrightarrow M_f, \quad
(a_1, \ldots, a_i) \longmapsto \sum a_im_i
$$
for some $f \in R$, $f \not \in \mathfrak p$. We may also assume that
$D(f) \subset U_i$. This means that $\wedge^i(M_f) = \wedge^i(M)_f$
is a flat $R_f$ module whose support is all of $\Spec(R_f)$.
By the above it is generated by a single element, namely
$m_1 \wedge \ldots \wedge m_i$. Hence $\wedge^i(M)_f \cong R_f/J$
for some pure ideal $J \subset R_f$ with $V(J) = \Spec(R_f)$.
Clearly this means that $J = (0)$, see
Lemma \ref{lemma-pure-ideal-determined-by-zero-set}.
Thus $m_1 \wedge \ldots \wedge m_i$ is a basis for
$\wedge^i(M_f)$ and it follows that the displayed map is injective
as well as surjective. This proves that $M$ is finite locally free
as desired.
\end{proof}







